Zusätzen wie dem Kurzreferat (Abstract), der Formulierung der Aufgabenstellung
oder einer Danksagung ist allen gemein, dass Sie mit einer Überschrift beginnen
und darunter einen knappen Text enthalten. Oft wäre es übertrieben hierfür ein
unnummieriertes Kapitel zu beginnen. Aus diesem Grund stellt die Vorlage die
\texttt{tucsimplesection}-Umgebung bereit. Ein Anwendungsbeispiel hierfür zeigt
Listing~\ref{lst:final:simplesection}.

\lstinputlisting[language={[LaTeX]{TeX}},style=block,escapechar=\!,float={htb},caption={Einfache Überschriften},label={lst:final:simplesection}]{code/final_simplesection.tex}

Abbildung~\ref{fig:final:simplesection} zeigt das Ergebnis, welches von der
Umgebung produziert wird. Der Abstand zwischen Überschrift und Inhalt kann
durch den optionalen Abstandparameter variiert werden. Standardmäßig ist
\verb+\bigskipamount+ voreingestellt.

\begin{figure}
\centering
\fbox{%
\begin{minipage}{0.95\linewidth}
\begin{tucsimplesection}{Überschrift}
\lipsum[1]
\end{tucsimplesection}
\end{minipage}}
\caption{Beispiel einer einfachen Überschrift}
\label{fig:final:simplesection}
\end{figure}

Der Autor ist selbst dafür verantwortlich, den Beginn einer neuen Seite zu
erzwingen (\verb+\cleardoublepage+). Die \verb+\vspace+-Kommandos dienen dazu
den Textblock samt Überschrift vertikal zu zentrieren. Der Abstand nach unten
wird dabei größer gewählt als nach oben. Dies ist optisch stimmiger.
